\documentclass[conference,10pt]{IEEEtran}
\usepackage{cite}
\usepackage{amsmath,amssymb,amsfonts}
\usepackage{graphicx}
\usepackage{xcolor}
\usepackage{booktabs}
\usepackage{url}
\usepackage[colorlinks=true,linkcolor=blue,citecolor=blue,urlcolor=blue]{hyperref}

% Every number in tables/figures is generated by `python -m inboxrl.plots` from results/*.json
% and runs/*/metrics.csv. \todo marks text that must be finalised from those generated files.
\newcommand{\todo}[1]{\textcolor{red}{[TODO: #1]}}
\newcommand{\inputorpending}[1]{\IfFileExists{#1}{\input{#1}}{\textit{(pending: run \texttt{python -m inboxrl.plots})}}}
\newcommand{\figorpending}[2]{\IfFileExists{figures/#1.pdf}{\includegraphics[width=#2]{figures/#1.pdf}}{\fbox{\parbox{0.9\linewidth}{\centering\small pending: figures/\detokenize{#1}.pdf}}}}

\begin{document}

\title{InboxRL: Pure Reinforcement Learning (GRPO) for\\Calendar Event Extraction from Institute Emails}

\author{
\IEEEauthorblockN{Hemant Patel (B23CH1018)}
\IEEEauthorblockA{IIT Jodhpur\\ \href{mailto:b23ch1018@iitj.ac.in}{b23ch1018@iitj.ac.in}}
\and
\IEEEauthorblockN{Kartik Gehlot (B23EE1088)}
\IEEEauthorblockA{IIT Jodhpur\\ \href{mailto:b23ee1088@iitj.ac.in}{b23ee1088@iitj.ac.in}}
}

\maketitle

\noindent\textbf{Code:} \url{https://github.com/hhemantpatel/inboxrl}

\begin{abstract}
Deadlines, quizzes and talks reach students as free-form institute emails. We train a small language model
(Qwen2.5-0.5B-Instruct with LoRA) to turn such an email into a JSON list of calendar events
\emph{purely with reinforcement learning}: no supervised fine-tuning, and labels are used only inside a verifiable
reward. Generation is cast as a token-level MDP and optimised with Group Relative Policy Optimization (GRPO).
The reward combines a format gate with a Hungarian bipartite match between predicted and gold events, and we
report two reward-hacking exploits we found in our first reward design and how we closed them.
On a template-disjoint synthetic test set and on held-out anonymised real emails, \todo{one-sentence result from results\_table.tex}.
\end{abstract}

\section{Problem}
A student receives many emails per day; missing one deadline is costly. Extraction must resolve relative dates
(``next Friday'' relative to the received date), ignore superseded dates (``postponed from 3 Nov''), handle
multiple events and noisy forwarded or Hinglish text, and emit machine-readable output that can be exported to an
\texttt{.ics} calendar. A 0.5B model runs locally (privacy) but, prompted zero-shot, often breaks the schema or
mis-resolves dates. The course rules forbid supervised methods, so we ask: \emph{can RL with a verifiable
reward alone teach this behaviour?}

\section{RL Formulation}
\textbf{State.} $s_t=(x_{\text{sys}},x_{\text{email}},d_0,y_{<t})$: system instructions, sender, subject and body
(truncated so the prompt is $\le$640 tokens), the received date $d_0$, and the tokens generated so far. The policy
embeds $s_t$ as a sequence of 896-dimensional hidden states (context $\le 640+192$ tokens).
\textbf{Action.} $a_t\in\mathcal V$, the next token; $|\mathcal V|=151{,}936$.
\textbf{Transition.} Deterministic, $s_{t+1}=s_t\circ a_t$; the episode ends at \texttt{<|im\_end|>} or after
$T=192$ tokens. \textbf{Discount} $\gamma=1$; the reward is sparse and terminal ($r_t=0$ for $t<T$).
\textbf{Policy} $\pi_\theta$: Qwen2.5-0.5B-Instruct \cite{qwen2024qwen25} with LoRA ($r{=}16,\alpha{=}32$, all
projections) \cite{hu2021lora}; the frozen base model is $\pi_{\text{ref}}$.

\subsection{Reward}
Let $P$ be the parsed predicted events and $G$ the gold events. If the output is not a JSON list that validates
against the schema, $r_{\text{fmt}}{=}-1$ and no match credit is given; otherwise $r_{\text{fmt}}{=}+1$.
For a pair $(p,g)$ with $\delta=\mathbb I[p.\text{date}=g.\text{date}]$:
\begin{equation}
S(p,g)=0.35\,\delta+\big(\delta+0.25(1-\delta)\big)\big(0.20\,s_{\text{time}}+0.25\,F_1^{\text{title}}+0.10\,s_{\text{type}}+0.10\,s_{\text{loc}}\big)
\end{equation}
where $F_1^{\text{title}}$ is bag-of-tokens F1 and $s_{\text{loc}}$ a fuzzy string ratio. The Hungarian algorithm
\cite{kuhn1955hungarian} finds the assignment $M$ maximising total similarity, and
\begin{equation}
r_{\text{match}}=\frac{\sum_{(p,g)\in M}S(p,g)}{\max(|P|,|G|)},
\end{equation}
so both missed and hallucinated events cost reward. $r_{\text{empty}}$ is $+1$ for $[\,]$ on a no-event email,
$-1$ for hallucinating on one, and $-0.5$ for answering $[\,]$ when events exist. $r_{\text{len}}{=}-1$ if any text
surrounds the JSON ($-0.5$ more beyond 800 characters). The total is
\begin{equation}
R=0.5\,r_{\text{fmt}}+2\,r_{\text{match}}+r_{\text{empty}}+0.25\,r_{\text{len}}\in[-1.75,\,2.5].
\end{equation}

\textbf{Reward hacking found and fixed.} Probing our first reward with hand-written policies exposed two exploits.
(i) Titles were scored with \texttt{token\_set\_ratio} plus a substring bonus, so the one-letter title ``a''
earned 2.45 of the maximum 2.5. (ii) The date was just one additive field, so an event on the \emph{wrong day}
kept 1.8/2.5. With token-F1 titles and date gating (Eq.~1) these fall to 2.0 and 0.83, while a wrong-date answer
still beats giving up ($[\,]$ scores 0), which keeps early learning signal. Unit tests pin these values and check that constant-$[\,]$ and ``copy every date'' policies score below an honest policy.

\begin{figure*}[t]\centering
\figorpending{comparison_real_test}{0.49\linewidth}\hfill
\figorpending{training_diagnostics}{0.49\linewidth}
\caption{Left: baselines vs.\ InboxRL on real emails. Right: KL, completion length and fraction of zero-variance groups during training.}
\end{figure*}

\section{Environment and Data}
\textbf{Synthetic generator} (our own code). Templates produce emails together with exact labels at five levels used as a curriculum:
L1 explicit dates; L2 varied formats (05/11/2026, Nov 5th, 5\,pm, EOD, noon); L3 relative dates; L4 postponements,
quoted reply chains and 2--4-event digests; L5 forwarded, noisy and partly Hinglish emails (``kal shaam 5 baje
LHC-101 mein quiz hai''). 15\% of emails contain no event. A quarter of all templates, plus two courses and two
venues, appear \emph{only} in the test split, so the 300 synthetic test emails measure generalisation to unseen phrasings rather than template memorisation (train 2000, val 200).
\textbf{Real emails.} \todo{N} institute emails from a Google Takeout export were anonymised (roll numbers, phones,
personal addresses, names) and hand-labelled; 60 form a held-out test set and the rest are mixed into the second
curriculum stage at 10\%. \texttt{EmailExtractionEnv} implements the MDP with \texttt{reset}/\texttt{step}.

\section{Algorithm}
GRPO \cite{shao2024deepseekmath} drops PPO's value network, which matters on a 4\,GB GPU. For each prompt $x$ we
sample a group of $G{=}4$ completions, and the advantage of completion $i$ is
$\hat A_i=(R_i-\operatorname{mean}_j R_j)/(\operatorname{std}_j R_j+\epsilon)$, shared by all its tokens.
With $\rho_{i,t}=\pi_\theta(y_{i,t}|x,y_{i,<t})/\pi_{\theta_{\text{old}}}(y_{i,t}|x,y_{i,<t})$ we maximise the
token-normalised clipped objective (TRL \cite{vonwerra2022trl}, \texttt{loss\_type=dapo})
\begin{equation}
\mathcal J(\theta)=\frac{1}{\sum_i|y_i|}\sum_{i}\sum_{t}\Big[\min\big(\rho_{i,t}\hat A_i,\operatorname{clip}(\rho_{i,t},1{\pm}0.2)\hat A_i\big)-\beta\,\mathbb D_{\text{KL}}^{(t)}\Big]
\end{equation}
with $\beta{=}0.04$ against $\pi_{\text{ref}}$. Each step uses 4 prompts $\times$ 4 completions, temperature 0.9,
learning rate $2{\times}10^{-5}$ (cosine, 50 warm-up steps), \todo{N} steps. \textbf{Curriculum}: the first
250 steps draw only L1--L2 emails, after which all levels are mixed.

\section{Experiments}
\textbf{Baselines:} constant $[\,]$ and random-guess policies (reward sanity checks), a regex+\texttt{dateparser}
rule system, and zero-/3-shot prompting of the same 0.5B model and of Qwen2.5-1.5B. No SFT baseline is included because the assignment forbids supervised training.
Event F1 counts a predicted event as correct when its date is right and its title token-F1 is $\ge$0.5; brackets
in the plots are 95\% bootstrap confidence intervals over emails.

\begin{table}[t]
\centering\caption{Test results (\%). GRPO: mean$\pm$std over 3 seeds.}\label{tab:results}
\resizebox{\linewidth}{!}{\inputorpending{results_table.tex}}
\end{table}

\begin{figure}[t]\centering
\figorpending{reward_vs_step}{\linewidth}
\caption{Training reward of GRPO rollouts and validation reward.}\label{fig:reward}
\end{figure}


\begin{table}[t]
\centering\caption{Ablations on the synthetic test set.}\label{tab:ablation}
\resizebox{0.9\linewidth}{!}{\inputorpending{ablation_table.tex}}
\end{table}

\textbf{Results.} \todo{Describe Table~\ref{tab:results} honestly: does GRPO beat zero-shot / few-shot / the 1.5B
model, by how much, and are the CIs separated? Mention per-level F1 (figures/levels\_*.pdf).}

\textbf{Ablations.} \todo{Summarise Table~\ref{tab:ablation}: curriculum, KL, format-only reward, group size.}

\textbf{Failure analysis.} \todo{2--3 real failure cases from results/completions/*.jsonl.}

\section{Conclusion}
\todo{Two sentences from the final results.} The trained adapter ships with a command-line tool that
turns an mbox export into an \texttt{.ics} file for Google Calendar.

\bibliographystyle{IEEEtran}
\bibliography{refs}

\end{document}
